\section{Training}\label{sec:training}

\para{Data and pseudo-labelling.} \ultra{} is trained on
LibriSpeech~\cite{panayotov2015librispeech} 960\,h combined with Libri-Light~\cite{kahn2020librilight}
pseudo-labelled by the pretrained streaming Zipformer INT8 model (R0, also the Zipformer comparator
in \cref{fig:money,tab:measured}), so its pseudo-labels inherit the Zipformer teacher's domain and
text style. Training logged \CandTrainKh\,h of audio: LibriSpeech until live WER saturates, then
Libri-Light (small, medium, large) and an EMA cooldown. In our runs LibriSpeech alone plateaus while new
audio keeps lowering WER, so for a model this small the lever is data diversity, not epochs.

\para{Recipe.} We optimise the joint RNN-T and auxiliary CTC loss (CTC weight \CtcWeight) with
AdamW, cosine learning-rate decay, SpecAugment~\cite{park2019specaugment}, speed perturbation, and
a small emission delay penalty (\DelayPenalty) to keep live-stream word latency low. The RNN-T loss
evaluates the joiner in time chunks without materialising the full $(B,T,U{+}1,V{+}1)$ lattice.
Training finishes with an exponential moving average (EMA)~\cite{polyak1992ema} cooldown phase,
yielding a single checkpoint whose weights are exported to prepared per-channel INT8 with zero
quantisation-aware retraining. The 5-gram language model (\texttt{M5.ulm}) used in the
utterance-segmented commit pass is trained on the standard LibriSpeech OpenSLR-11 text corpus.
